\documentclass[11pt]{article}
\usepackage{fontspec}
\setmainfont{Times New Roman}
\setsansfont{Arial}
\setmonofont{Consolas}
\usepackage{amsmath,amssymb}
\usepackage{geometry}
\usepackage{microtype}
\geometry{a4paper,margin=3cm}
\setlength{\parindent}{1.25em}
\setlength{\parskip}{0pt}
\emergencystretch=30em
\allowdisplaybreaks
\raggedbottom

\begin{document}

% Унит: Папер 21. Дирецт Интерславиц Латин транслатион фром те Герман финал-аудитед соурце слице.

\section*{21. Формалны вариационны рачун и диференциалне инварианты}
\begin{center}
Encyklopädie d. math. Wiss. II, 3 (1922), стр. 68--71 (в: R. Weitzenböck, Differentialinvarianten)
\end{center}

\setcounter{footnote}{148}
\subsection*{28. Формалны вариационны рачун и диференциалне инварианты.\footnote{Та точка походи од E. Noether.}}

Образованье всих диференциалных инвариантов и извод главных теоремов можно најједнообразнеје провести формалными методами вариационного рачуна. То начело иде назад к Риман\textsuperscript{72)} и своје главно формо-теоретично развитје добило јест у Липшиц\textsuperscript{73)}.

Ако выходимо из вариационној задачи, ктора принадлежи хомогенној форме $f(dx)$, кде
\[
\left|\frac{\partial^2 f}{\partial dx_i\,\partial dx_k}\right|\ne0
\]
јест предположено, именно
\begin{equation}
 \delta\int_{t_0}^{t_1} f(x')\,dt=0,
 \qquad \left(x_i'=\frac{dx_i}{dt}\right),
\tag{140}
\end{equation}
тогда интеграција по частјах води к инвариантној системи једначень
\[
 2\psi_i=\frac{d}{dt}\frac{\partial f}{\partial x_i'}-\frac{\partial f}{\partial x_i}=0,
\]
кде Лагрангеове изразы $\psi_i$ сут контрагредиентне к $dx$. Формално то најпростеје види се, ако $\psi_i$ определити безинтегралноју формалноју идентичностју, ктора соответује интеграцији по частјах (централно једначенье Лагрангеа\footnote{Названье за специално $f$ походи од Хеун; гл. јего чланок IV, 1 11, p. 447.})
\[
 \delta f-df_\delta=-2\sum\psi_i(d,d)\,\delta x_i,
 \qquad
 \left(f_\delta=\sum\frac{\partial f}{\partial dx_i}\,\delta x_i\right),
\]
кде и $f$ и $df_\delta$ сут инварианты, и зато такоже права страна. Особливо за квадратну форму добива се
\[
 \delta\sum g_{ik}dx_i dx_k-2d\sum g_{ik}dx_i\delta x_k
 =-2\sum\psi_\mu(d,d)\,\delta x_\mu.
\]
Ако тут заменити $d$ когредиентным $(d+\lambda\delta)$ и сравнити степены по $\lambda$, добива се соответна система идентичностиј:
\begin{equation}
\left\{
\begin{aligned}
D\sum g_{ik}dx_i dx_k-2d\sum g_{ik}dx_iDx_k
   &=-2\sum\psi_\mu(d,d)\,Dx_\mu,\\
D\sum g_{ik}dx_i\delta x_k-\delta\sum g_{ik}dx_iDx_k
   -d\sum g_{ik}\delta x_iDx_k
   &=-2\sum\psi_\mu(d,\delta)\,Dx_\mu,\\
D\sum g_{ik}\delta x_i\delta x_k-2\delta\sum g_{ik}\delta x_iDx_k
   &=-2\sum\psi_\mu(\delta,\delta)\,Dx_\mu.
\end{aligned}
\right.
\tag{141}
\end{equation}
Из того $\psi_\mu(d,\delta)$, и особливо такоже $\psi_\mu(d,d)$ и $\psi_\mu(\delta,\delta)$, показывајут се како контрагредиентне векторы. Из (141) израчунава се
\[
 \psi_\mu(d,\delta)=\sum g_{i\mu}\,d\delta x_i+
 \sum\left[\begin{smallmatrix}ik\\ \mu\end{smallmatrix}\right]dx_i\delta x_k,
\]
и из того добива се когредиентны вектор
\begin{equation}
 p^\sigma(d,\delta)=\sum g^{\sigma\mu}\psi_\mu
 =d\delta x_\sigma+
 \sum\left\{\begin{smallmatrix}ik\\ \sigma\end{smallmatrix}\right\}dx_i\delta x_k.
\tag{142}
\end{equation}
Ако положити $p^\nu(d,d)=0$, то очивидно даје једначеньа геодезичных линиј.

Знанье когредиентного вектора $p^\sigma(d,\delta)$ веде непосреднье к ковариантној деривацији (ср. нр. 19). Ако напр. $h(d,\delta)$ јест форма двох редов $dx$ и $\delta x$, тогда израз
\begin{equation}
 h^{(1)}(dx,\delta x)=\delta h-
 \sum\frac{\partial h}{\partial dx_\sigma}p^\sigma(d,\delta)-
 \sum\frac{\partial h}{\partial \delta x_\sigma}p^\sigma(\delta,\delta)
\tag{143}
\end{equation}
како разлика инвариантов сам јест инвариант, и јест образованы тако, же друге диференциалы се одстраньујут. Зато $h^{(1)}(dx,dx)$ представја форму само с првыми диференциалами, ковариантну деривацију од $h$. Соответно важи, ако $h$ содржи выше редов $d^{(1)}x,\ldots,d^{(\lambda)}x$.\footnote{Липшиц, J. f. Math. 72 (1870), p. 17.}

На том самом начелу основано јест образованье формы кривины у Риман и Липшиц. Преходи се од $\delta^2 f$ к ``нормалној форме другеј вариације''
\begin{equation}
 \Omega=\delta^2\sum g_{ik}dx_i dx_k
 -2d\delta\sum g_{ik}dx_i\delta x_k
 +d^2\sum g_{ik}\delta x_i\delta x_k,
\tag{144}
\end{equation}
ктора, како агрегат инвариантов, сама знову јест инвариант, и в кторој тепер, в обобченьу централного једначеньа Лагрангеа, третје диференциалы сут елиминоване. Елиминација другых диференциалов дела се в аналогији с ковариантным диференцираньем чрез образованье разлике
\begin{equation}
 K=\Omega-2\sum\{p^\sigma(d,\delta)\psi_\sigma(d,\delta)-p^\sigma(\delta,\delta)\psi_\sigma(d,d)\},
\tag{145}
\end{equation}
што можно такоже писати како\footnote{Липшиц, J. f. Math. 82 (1876), § 1.}
\begin{equation}
\begin{aligned}
 \frac12 K={}&d\sum\psi_\sigma(d,\delta)\delta x_\sigma
 -\delta\sum\psi_\sigma(d,d)\delta x_\sigma\\
 &-\sum\{p^\sigma(d,\delta)\psi_\sigma(d,\delta)-p^\sigma(\delta,\delta)\psi_\sigma(d,d)\}.
\end{aligned}
\tag{146}
\end{equation}
Закон образованьа $K$ можно изректи и тако, же в $\Omega$ друге диференциалы елиминујут се чрез положанье $p(d,\delta)$ и $p(\delta,\delta)$, соответно правых стран (141), равными нули. То јест определьенье формы кривины по Риман (ср. нр. 21), при чем треба изрично приметити -- што (146) особливо јасно показује -- же то положанье к нули може наступити само после проведеного диференцираньа; формалны начин образованьа у Липшиц обходи ту трудност. Соответно можно бы было определити и ковариантну деривацију $h^{(1)}$ (143) тако, же в $\delta h$ друге диференциалы елиминоване сут чрез положанье $p(d,\delta),\ldots$ равными нули.

Из формы кривины чрез сукцесивно образованье ковариантных деривациј добива се ред форм $[\Phi_4,\Phi_5,\ldots$ у Кристофел], кторых пројективне инварианты, после доданьа $f$, по Кристофел и Ричи изчерпавајут целост диференциалных инвариантов квадратној формы, а их еквивалентност взгледом линеарных трансформациј без дального условја интеграбилности изрека еквивалентност двох квадратных диференциалных форм (ср. нр. 22). Внутрны основ того теорема редукције опираје се на ексистенцији Римановых нормалных координат, кторе наставајут из вариационној задачи (140) (ср. нр. 20): оне трансформујут геодезичне линије идуче чрез точку в праве линије, и зато при либовольној трансформацији од $x$ трансформујут се линеарно. Образованье всих инвариантов и питанье еквивалентности сут зато сведене к соответному питаньу за поједине хомогене чести в развоју од $f$ по нормалных координатах, и показује се, же те чести линеарно составльајут се из $f,\Phi_4,\Phi_5,\ldots$. За квадратне формы то јест в подробностјах проведено у Вермеил\textsuperscript{92)}, в связи с нотоју E. Noether\textsuperscript{93)}. По тој ноте теоремы и методы оставајут валидне, ако в основу положити либовольны диференциалны израз (ср. нр. 22 и 23); тым показује се досег методов формалного вариационного рачуна в противности к методам чистеј теорији елиминације.\footnote{Ср. к тым изложеньам и за дальне геометричне толкованьа семинарне предаваньа Ф. Клеин, наведене под ``Литература''.}

Положанье $p(d,\delta)$ равным нули Леви-Чивита\textsuperscript{134)} геометрично изтолковал јест како ``паралелно пренешенье вектора'' (ср. такоже Хессенберг\textsuperscript{133)}). То поньатје, аксиоматично ухвачено од Вејл, в сушчности јест ограничено на квадратне формы, але допушча обобченье другого рода. Оно оставаје при разширеньу групы (множенье $g_{ik}$ либовольноју функцијеју), скоро кроме квадратној формы такоже линеарна форма положена јест в основу. Тогда $p(d,\delta)$ јест једнозначно определьен, и можно проводити образованьа инвариантов соответные горньим и определьати геодезичне координаты; но $p(d,d)=0$ тепер выше не наставаје из вариационној задачи!\footnote{Х. Вејл, Reine Infinitesimalgeometrie, Math. Ztschr. 2 (1918), p. 384--411, и Raum, Zeit, Materie (3. и 4. изд.).} Питаньа о теорему редукције и еквивалентности тут јешче не сут обче разматриване; само за инварианты најнижшего реда R. Weitzenböck дал јест теорем редукције.\footnote{Wien. Ber. 129 (1920), p. 683 и 697.}

Наконец треба указати на обче ``инвариантне вариацијске проблемы'', т. ј. на таке, при кторых просты или многократны интеграл $J$ јест инвариантен взходно к некој групе в смыслу Lie. Специалными физично интересными случајами занимајут се работы Хамел\footnote{Math. Ann. 59 (1904), p. 416--434; Ztschr. Math. Phys. 50 (1904), p. 1--54.}, Херглотз\footnote{Ann. d. Phys. (4) 36 (1911), p. 511, особливо § 9.}, Лорентз\footnote{Verslag der Amsterd. Akad., април, септ. 1916, окт. 1916, мај 1917.}, Фоккер\footnote{Тамже, јан. 1917.}, Вејл\footnote{Ann. d. Phys. (4) 54 (1917), p. 117, § 2.} и Клеин\footnote{Gött. Nachr. 1917, p. 469--482; тамже 1918, p. 171--189.}. Принципиална формулација у E. Noether\footnote{Gött. Nachr. 1918, p. 235--257.} показује, же инвариантности $J$ взходно к $G_\rho$ (конечна група с $\rho$ сушчствеными параметрами) соответујут $\rho$ линеарно независне дивергенције; инвариантности взходно к бесконечној групе, ктора содржи $\rho$ либовольных функциј до $\sigma$-теј деривације, соответујут $\rho$ зависимости меджу Лагрангеовыми изразами и их деривацијами до $\sigma$-того реда. В обојих случајах важи обрат. Понеже Лагрангеове изразы ставајут релативными инвариантами групы, имајемо вместе с тым процес, кторы образује инварианты.

\begin{center}
(Закльучено в марцу 1921.)
\end{center}

\end{document}

